\documentclass[12pt,a4paper]{article}
\usepackage[utf8]{vietnam}
\usepackage{amsmath,amssymb}
\usepackage{tabularx}
\usepackage{tikz}
\usepackage{xcolor}
\usepackage[top=0.6cm, bottom=0.6cm, left=0.6cm, right=0.6cm]{geometry}
\usepackage[dethi]{ex_test}

\SAOPTN{widthSA=3.2,heightSA=0.62}

\setlength{\parskip}{0pt}

% Dòng kẻ chấm mờ đều đặn chuẩn viết tay
\newcommand{\dongke}[1]{%
  \begingroup%
  \setlength{\parskip}{0pt}%
  \baselineskip=18.5pt%
  \vspace{1.0mm}%
  \foreach \i in {1,...,#1}{%
    \noindent{\color{gray!60}\dotfill}\par%
  }%
  \vspace{0.6mm}%
  \endgroup%
}

% Định nghĩa ô bài tập tự luận trong bảng
\newcommand{\oBaiTap}[3][6]{%
  \vspace{1.0mm}%
  \noindent\textbf{#2} #3\par%
  \dongke{#1}%
}

\begin{document}
\pagestyle{empty}
% =============================================================================
% ========================= TOÁN 7 - TUẦN 06 - BUỔI 01: MÃ ĐỀ 101 =============
% ========================= MẶT TRƯỚC (TRANG 1/2) ==============================
% =============================================================================

\noindent
\begin{tabularx}{\linewidth}{|X|c|}
\hline
\textbf{PHIẾU KIỂM TRA ĐỊNH KỲ TOÁN 7 - TUẦN 06 - BUỔI 01} & \textbf{Mã đề: 101 - Trang 1/2} \\
\hline
\end{tabularx}

\vspace{1.0mm}

\noindent
\begin{tabularx}{\linewidth}{|p{3.2cm}|X|}
\hline
\centering\textbf{ĐIỂM SỐ} & \centering\textbf{LỜI NHẬN XÉT CỦA THẦY CÔ} \tabularnewline
\hline
\rule{0pt}{1.8cm} & \tabularnewline
\hline
\end{tabularx}

\vspace{1.2mm}
\noindent\textbf{Câu 1 (0,25 điểm).} Tìm căn bậc hai số học của $49$. \shortans{7}\par\vspace{0.8mm}
\noindent\textbf{Câu 2 (0,25 điểm).} Tính giá trị của biểu thức: $|-4{,}5| + |2{,}5|$. \shortans{7}\par\vspace{0.8mm}
\noindent\textbf{Câu 3 (0,25 điểm).} Làm tròn số $15{,}738$ đến chữ số thập phân thứ hai. \shortans{15{,}74}\par\vspace{0.8mm}
\noindent\textbf{Câu 4 (0,25 điểm).} Cho góc $\widehat{AOB} = 65^\circ$. Số đo của góc kề bù với góc $\widehat{AOB}$ là: \shortans{115^\circ}\par\vspace{0.8mm}
\noindent\textbf{Câu 5 (0,25 điểm).} Hai đường thẳng cắt nhau tạo thành một góc bằng $55^\circ$. Số đo góc đối đỉnh với nó là: \shortans{55^\circ}\par\vspace{0.8mm}
\noindent\textbf{Câu 6 (0,25 điểm).} Cho tia $Oz$ là tia phân giác của góc $\widehat{xOy} = 70^\circ$. Số đo góc $\widehat{xOz}$ là: \shortans{35^\circ}\par\vspace{0.8mm}
\noindent\textbf{Câu 7 (0,25 điểm).} Tính giá trị của biểu thức: $\left(-\dfrac{2}{3}\right)^2 + \dfrac{5}{9}$. \shortans{1}\par\vspace{0.8mm}
\noindent\textbf{Câu 8 (0,25 điểm).} Cho $a // b$, một đường thẳng cắt $a$ và $b$. Biết một góc so le trong bằng $60^\circ$, số đo góc so le trong còn lại là: \shortans{60^\circ}\par\vspace{0.8mm}
\noindent\textbf{Câu 9 (0,25 điểm).} Tìm số thực $x \ge 0$, biết $\sqrt{x} = 5$. \shortans{25}\par\vspace{0.8mm}
\noindent\textbf{Câu 10 (0,25 điểm).} Một chiếc cặp giá $200\,000$ đồng được giảm giá $10\%$. Số tiền được giảm là (đồng): \shortans{20\,000}\par\vspace{0.8mm}
\vspace{0.8mm}
\noindent\textbf{Câu 11 (2,0 điểm). Thực hiện phép tính (hợp lý nếu có thể):}
\vspace{0.6mm}

\noindent
\begin{tabularx}{\linewidth}{|X|X|}
\hline
\oBaiTap[9]{a)}{$\dfrac{3}{7} + \left(-\dfrac{5}{2}\right) + \dfrac{4}{7}$} & 
\oBaiTap[9]{b)}{$\dfrac{4}{9} \cdot \left(-\dfrac{5}{13}\right) + \dfrac{4}{9} \cdot \left(-\dfrac{8}{13}\right)$} \\
\hline
\oBaiTap[9]{c)}{$\left(-\dfrac{1}{2}\right)^3 : \dfrac{1}{4} + 2{,}5 \cdot (-2)$} & 
\oBaiTap[9]{d)}{$\sqrt{\dfrac{25}{4}} - |-1{,}5| + \sqrt{16}$} \\
\hline
\end{tabularx}

\newpage
% =============================================================================
% ========================= TOÁN 7 - TUẦN 06 - BUỔI 01: MÃ ĐỀ 101 =============
% ========================= MẶT SAU (TRANG 2/2) ================================
% =============================================================================

\noindent
\begin{tabularx}{\linewidth}{|X|c|}
\hline
\textbf{TOÁN 7 - TUẦN 06 - BUỔI 01} & \textbf{Mã đề: 101 - Trang 2/2} \\
\hline
\end{tabularx}

\vspace{1.2mm}

\noindent\textbf{Câu 12 (2,0 điểm). Tìm số hữu tỉ $x$, biết:}
\vspace{0.6mm}

\noindent
\begin{tabularx}{\linewidth}{|X|X|}
\hline
\oBaiTap[5]{a)}{$x + \dfrac{2}{3} = \dfrac{5}{6}$} & 
\oBaiTap[5]{b)}{$\dfrac{3}{4}x - \dfrac{1}{2} = \dfrac{1}{4}$} \\
\hline
\oBaiTap[5]{c)}{$\left(x - \dfrac{1}{2}\right)^2 = \dfrac{9}{16}$} & 
\oBaiTap[5]{d)}{$|2x - 1| = 3$} \\
\hline
\end{tabularx}

\vspace{1.2mm}

\noindent
\begin{tabularx}{\linewidth}{|X|}
\hline
\vspace{0.5mm}
\textbf{Câu 13 (2,5 điểm).} Cho hai góc kề bù $\widehat{xOy}$ và $\widehat{yOz}$. Biết $\widehat{xOy} = 60^\circ$.\par
\textbf{a) (1,25 điểm)} Tính số đo góc $\widehat{yOz}$.\par
\textbf{b) (1,25 điểm)} Vẽ tia $Om$ là tia phân giác của góc $\widehat{yOz}$. Tính số đo góc $\widehat{xOm}$.\par
\vspace{0.4mm}
\noindent\textit{Bài giải:}\dongke{13}
\\ \hline
\end{tabularx}

\vspace{1.2mm}

\noindent
\begin{tabularx}{\linewidth}{|X|}
\hline
\vspace{0.5mm}
\textbf{Câu 14 (1,0 điểm).} Tìm giá trị nhỏ nhất của biểu thức $A = |x - 3| + 2026$. Chỉ ra giá trị của $x$ để biểu thức đạt giá trị nhỏ nhất đó.\par
\vspace{0.4mm}
\noindent\textit{Bài giải:}\dongke{8}
\\ \hline
\end{tabularx}

\newpage
% =============================================================================
% ========================= TOÁN 7 - TUẦN 06 - BUỔI 01: MÃ ĐỀ 102 =============
% ========================= MẶT TRƯỚC (TRANG 1/2) ==============================
% =============================================================================

\noindent
\begin{tabularx}{\linewidth}{|X|c|}
\hline
\textbf{PHIẾU KIỂM TRA ĐỊNH KỲ TOÁN 7 - TUẦN 06 - BUỔI 01} & \textbf{Mã đề: 102 - Trang 1/2} \\
\hline
\end{tabularx}

\vspace{1.0mm}

\noindent
\begin{tabularx}{\linewidth}{|p{3.2cm}|X|}
\hline
\centering\textbf{ĐIỂM SỐ} & \centering\textbf{LỜI NHẬN XÉT CỦA THẦY CÔ} \tabularnewline
\hline
\rule{0pt}{1.8cm} & \tabularnewline
\hline
\end{tabularx}

\vspace{1.2mm}
\noindent\textbf{Câu 1 (0,25 điểm).} Tìm căn bậc hai số học của $81$. \shortans{9}\par\vspace{0.8mm}
\noindent\textbf{Câu 2 (0,25 điểm).} Tính giá trị của biểu thức: $|-3{,}8| + |2{,}2|$. \shortans{6}\par\vspace{0.8mm}
\noindent\textbf{Câu 3 (0,25 điểm).} Làm tròn số $24{,}816$ đến chữ số thập phân thứ hai. \shortans{24{,}82}\par\vspace{0.8mm}
\noindent\textbf{Câu 4 (0,25 điểm).} Cho góc $\widehat{AOB} = 75^\circ$. Số đo của góc kề bù với góc $\widehat{AOB}$ là: \shortans{105^\circ}\par\vspace{0.8mm}
\noindent\textbf{Câu 5 (0,25 điểm).} Hai đường thẳng cắt nhau tạo thành một góc bằng $68^\circ$. Số đo góc đối đỉnh với nó là: \shortans{68^\circ}\par\vspace{0.8mm}
\noindent\textbf{Câu 6 (0,25 điểm).} Cho tia $Oz$ là tia phân giác của góc $\widehat{xOy} = 84^\circ$. Số đo góc $\widehat{xOz}$ là: \shortans{42^\circ}\par\vspace{0.8mm}
\noindent\textbf{Câu 7 (0,25 điểm).} Tính giá trị của biểu thức: $\left(-\dfrac{3}{4}\right)^2 + \dfrac{7}{16}$. \shortans{1}\par\vspace{0.8mm}
\noindent\textbf{Câu 8 (0,25 điểm).} Cho $a // b$, biết một góc đồng vị có số đo bằng $70^\circ$, số đo góc đồng vị tương ứng là: \shortans{70^\circ}\par\vspace{0.8mm}
\noindent\textbf{Câu 9 (0,25 điểm).} Tìm số thực $x \ge 0$, biết $\sqrt{x} = 7$. \shortans{49}\par\vspace{0.8mm}
\noindent\textbf{Câu 10 (0,25 điểm).} Một chiếc áo giá $300\,000$ đồng được giảm giá $10\%$. Số tiền được giảm là (đồng): \shortans{30\,000}\par\vspace{0.8mm}
\vspace{0.8mm}
\noindent\textbf{Câu 11 (2,0 điểm). Thực hiện phép tính (hợp lý nếu có thể):}
\vspace{0.6mm}

\noindent
\begin{tabularx}{\linewidth}{|X|X|}
\hline
\oBaiTap[9]{a)}{$\dfrac{2}{9} + \left(-\dfrac{7}{3}\right) + \dfrac{7}{9}$} & 
\oBaiTap[9]{b)}{$\dfrac{5}{8} \cdot \left(-\dfrac{7}{15}\right) + \dfrac{5}{8} \cdot \left(-\dfrac{8}{15}\right)$} \\
\hline
\oBaiTap[9]{c)}{$\left(-\dfrac{1}{3}\right)^2 : \dfrac{1}{9} - 1{,}5 \cdot 4$} & 
\oBaiTap[9]{d)}{$\sqrt{\dfrac{49}{4}} - |-2{,}5| + \sqrt{25}$} \\
\hline
\end{tabularx}

\newpage
% =============================================================================
% ========================= TOÁN 7 - TUẦN 06 - BUỔI 01: MÃ ĐỀ 102 =============
% ========================= MẶT SAU (TRANG 2/2) ================================
% =============================================================================

\noindent
\begin{tabularx}{\linewidth}{|X|c|}
\hline
\textbf{TOÁN 7 - TUẦN 06 - BUỔI 01} & \textbf{Mã đề: 102 - Trang 2/2} \\
\hline
\end{tabularx}

\vspace{1.2mm}

\noindent\textbf{Câu 12 (2,0 điểm). Tìm số hữu tỉ $x$, biết:}
\vspace{0.6mm}

\noindent
\begin{tabularx}{\linewidth}{|X|X|}
\hline
\oBaiTap[5]{a)}{$x + \dfrac{3}{4} = \dfrac{7}{8}$} & 
\oBaiTap[5]{b)}{$\dfrac{2}{3}x - \dfrac{1}{6} = \dfrac{1}{2}$} \\
\hline
\oBaiTap[5]{c)}{$\left(x - \dfrac{1}{3}\right)^2 = \dfrac{4}{9}$} & 
\oBaiTap[5]{d)}{$|2x - 3| = 5$} \\
\hline
\end{tabularx}

\vspace{1.2mm}

\noindent
\begin{tabularx}{\linewidth}{|X|}
\hline
\vspace{0.5mm}
\textbf{Câu 13 (2,5 điểm).} Cho hai góc kề bù $\widehat{xOy}$ và $\widehat{yOz}$. Biết $\widehat{xOy} = 70^\circ$.\par
\textbf{a) (1,25 điểm)} Tính số đo góc $\widehat{yOz}$.\par
\textbf{b) (1,25 điểm)} Vẽ tia $Om$ là tia phân giác của góc $\widehat{yOz}$. Tính số đo góc $\widehat{xOm}$.\par
\vspace{0.4mm}
\noindent\textit{Bài giải:}\dongke{13}
\\ \hline
\end{tabularx}

\vspace{1.2mm}

\noindent
\begin{tabularx}{\linewidth}{|X|}
\hline
\vspace{0.5mm}
\textbf{Câu 14 (1,0 điểm).} Tìm giá trị nhỏ nhất của biểu thức $A = |x - 5| + 2025$. Chỉ ra giá trị của $x$ để biểu thức đạt giá trị nhỏ nhất đó.\par
\vspace{0.4mm}
\noindent\textit{Bài giải:}\dongke{8}
\\ \hline
\end{tabularx}

\newpage
% =============================================================================
% ========================= TOÁN 7 - TUẦN 06 - BUỔI 01: MÃ ĐỀ 103 =============
% ========================= MẶT TRƯỚC (TRANG 1/2) ==============================
% =============================================================================

\noindent
\begin{tabularx}{\linewidth}{|X|c|}
\hline
\textbf{PHIẾU KIỂM TRA ĐỊNH KỲ TOÁN 7 - TUẦN 06 - BUỔI 01} & \textbf{Mã đề: 103 - Trang 1/2} \\
\hline
\end{tabularx}

\vspace{1.0mm}

\noindent
\begin{tabularx}{\linewidth}{|p{3.2cm}|X|}
\hline
\centering\textbf{ĐIỂM SỐ} & \centering\textbf{LỜI NHẬN XÉT CỦA THẦY CÔ} \tabularnewline
\hline
\rule{0pt}{1.8cm} & \tabularnewline
\hline
\end{tabularx}

\vspace{1.2mm}
\noindent\textbf{Câu 1 (0,25 điểm).} Tìm căn bậc hai số học của $64$. \shortans{8}\par\vspace{0.8mm}
\noindent\textbf{Câu 2 (0,25 điểm).} Tính giá trị của biểu thức: $|-5{,}4| + |2{,}6|$. \shortans{8}\par\vspace{0.8mm}
\noindent\textbf{Câu 3 (0,25 điểm).} Làm tròn số $18{,}364$ đến chữ số thập phân thứ hai. \shortans{18{,}36}\par\vspace{0.8mm}
\noindent\textbf{Câu 4 (0,25 điểm).} Cho góc $\widehat{AOB} = 50^\circ$. Số đo của góc kề bù với góc $\widehat{AOB}$ là: \shortans{130^\circ}\par\vspace{0.8mm}
\noindent\textbf{Câu 5 (0,25 điểm).} Hai đường thẳng cắt nhau tạo thành một góc bằng $42^\circ$. Số đo góc đối đỉnh với nó là: \shortans{42^\circ}\par\vspace{0.8mm}
\noindent\textbf{Câu 6 (0,25 điểm).} Cho tia $Oz$ là tia phân giác của góc $\widehat{xOy} = 96^\circ$. Số đo góc $\widehat{xOz}$ là: \shortans{48^\circ}\par\vspace{0.8mm}
\noindent\textbf{Câu 7 (0,25 điểm).} Tính giá trị của biểu thức: $\left(-\dfrac{1}{2}\right)^3 + \dfrac{9}{8}$. \shortans{1}\par\vspace{0.8mm}
\noindent\textbf{Câu 8 (0,25 điểm).} Cho $a // b$, biết một góc so le trong bằng $55^\circ$, số đo góc so le trong còn lại là: \shortans{55^\circ}\par\vspace{0.8mm}
\noindent\textbf{Câu 9 (0,25 điểm).} Tìm số thực $x \ge 0$, biết $\sqrt{x} = 8$. \shortans{64}\par\vspace{0.8mm}
\noindent\textbf{Câu 10 (0,25 điểm).} Một đôi giày giá $400\,000$ đồng được giảm giá $10\%$. Số tiền được giảm là (đồng): \shortans{40\,000}\par\vspace{0.8mm}
\vspace{0.8mm}
\noindent\textbf{Câu 11 (2,0 điểm). Thực hiện phép tính (hợp lý nếu có thể):}
\vspace{0.6mm}

\noindent
\begin{tabularx}{\linewidth}{|X|X|}
\hline
\oBaiTap[9]{a)}{$\dfrac{4}{11} + \left(-\dfrac{5}{2}\right) + \dfrac{7}{11}$} & 
\oBaiTap[9]{b)}{$\dfrac{3}{7} \cdot \left(-\dfrac{9}{17}\right) + \dfrac{3}{7} \cdot \left(-\dfrac{8}{17}\right)$} \\
\hline
\oBaiTap[9]{c)}{$\left(-\dfrac{1}{2}\right)^2 \cdot 8 - 3{,}5 \cdot 2$} & 
\oBaiTap[9]{d)}{$\sqrt{\dfrac{81}{4}} - |-3{,}5| + \sqrt{36}$} \\
\hline
\end{tabularx}

\newpage
% =============================================================================
% ========================= TOÁN 7 - TUẦN 06 - BUỔI 01: MÃ ĐỀ 103 =============
% ========================= MẶT SAU (TRANG 2/2) ================================
% =============================================================================

\noindent
\begin{tabularx}{\linewidth}{|X|c|}
\hline
\textbf{TOÁN 7 - TUẦN 06 - BUỔI 01} & \textbf{Mã đề: 103 - Trang 2/2} \\
\hline
\end{tabularx}

\vspace{1.2mm}

\noindent\textbf{Câu 12 (2,0 điểm). Tìm số hữu tỉ $x$, biết:}
\vspace{0.6mm}

\noindent
\begin{tabularx}{\linewidth}{|X|X|}
\hline
\oBaiTap[5]{a)}{$x - \dfrac{1}{3} = \dfrac{1}{6}$} & 
\oBaiTap[5]{b)}{$\dfrac{4}{5}x + \dfrac{1}{5} = 1$} \\
\hline
\oBaiTap[5]{c)}{$\left(x + \dfrac{1}{2}\right)^2 = \dfrac{25}{16}$} & 
\oBaiTap[5]{d)}{$|2x + 1| = 5$} \\
\hline
\end{tabularx}

\vspace{1.2mm}

\noindent
\begin{tabularx}{\linewidth}{|X|}
\hline
\vspace{0.5mm}
\textbf{Câu 13 (2,5 điểm).} Cho hai góc kề bù $\widehat{xOy}$ và $\widehat{yOz}$. Biết $\widehat{xOy} = 50^\circ$.\par
\textbf{a) (1,25 điểm)} Tính số đo góc $\widehat{yOz}$.\par
\textbf{b) (1,25 điểm)} Vẽ tia $Om$ là tia phân giác của góc $\widehat{yOz}$. Tính số đo góc $\widehat{xOm}$.\par
\vspace{0.4mm}
\noindent\textit{Bài giải:}\dongke{13}
\\ \hline
\end{tabularx}

\vspace{1.2mm}

\noindent
\begin{tabularx}{\linewidth}{|X|}
\hline
\vspace{0.5mm}
\textbf{Câu 14 (1,0 điểm).} Tìm giá trị nhỏ nhất của biểu thức $A = |x + 4| + 2026$. Chỉ ra giá trị của $x$ để biểu thức đạt giá trị nhỏ nhất đó.\par
\vspace{0.4mm}
\noindent\textit{Bài giải:}\dongke{8}
\\ \hline
\end{tabularx}

\newpage
% =============================================================================
% ========================= TOÁN 7 - TUẦN 06 - BUỔI 01: MÃ ĐỀ 104 =============
% ========================= MẶT TRƯỚC (TRANG 1/2) ==============================
% =============================================================================

\noindent
\begin{tabularx}{\linewidth}{|X|c|}
\hline
\textbf{PHIẾU KIỂM TRA ĐỊNH KỲ TOÁN 7 - TUẦN 06 - BUỔI 01} & \textbf{Mã đề: 104 - Trang 1/2} \\
\hline
\end{tabularx}

\vspace{1.0mm}

\noindent
\begin{tabularx}{\linewidth}{|p{3.2cm}|X|}
\hline
\centering\textbf{ĐIỂM SỐ} & \centering\textbf{LỜI NHẬN XÉT CỦA THẦY CÔ} \tabularnewline
\hline
\rule{0pt}{1.8cm} & \tabularnewline
\hline
\end{tabularx}

\vspace{1.2mm}
\noindent\textbf{Câu 1 (0,25 điểm).} Tìm căn bậc hai số học của $100$. \shortans{10}\par\vspace{0.8mm}
\noindent\textbf{Câu 2 (0,25 điểm).} Tính giá trị của biểu thức: $|-6{,}3| + |2{,}7|$. \shortans{9}\par\vspace{0.8mm}
\noindent\textbf{Câu 3 (0,25 điểm).} Làm tròn số $32{,}645$ đến chữ số thập phân thứ hai. \shortans{32{,}65}\par\vspace{0.8mm}
\noindent\textbf{Câu 4 (0,25 điểm).} Cho góc $\widehat{AOB} = 80^\circ$. Số đo của góc kề bù với góc $\widehat{AOB}$ là: \shortans{100^\circ}\par\vspace{0.8mm}
\noindent\textbf{Câu 5 (0,25 điểm).} Hai đường thẳng cắt nhau tạo thành một góc bằng $72^\circ$. Số đo góc đối đỉnh với nó là: \shortans{72^\circ}\par\vspace{0.8mm}
\noindent\textbf{Câu 6 (0,25 điểm).} Cho tia $Oz$ là tia phân giác của góc $\widehat{xOy} = 64^\circ$. Số đo góc $\widehat{xOz}$ là: \shortans{32^\circ}\par\vspace{0.8mm}
\noindent\textbf{Câu 7 (0,25 điểm).} Tính giá trị của biểu thức: $\left(-\dfrac{2}{5}\right)^2 + \dfrac{21}{25}$. \shortans{1}\par\vspace{0.8mm}
\noindent\textbf{Câu 8 (0,25 điểm).} Cho $a // b$, biết một góc đồng vị có số đo bằng $80^\circ$, số đo góc đồng vị tương ứng là: \shortans{80^\circ}\par\vspace{0.8mm}
\noindent\textbf{Câu 9 (0,25 điểm).} Tìm số thực $x \ge 0$, biết $\sqrt{x} = 6$. \shortans{36}\par\vspace{0.8mm}
\noindent\textbf{Câu 10 (0,25 điểm).} Một quyển sách giá $150\,000$ đồng được giảm giá $10\%$. Số tiền được giảm là (đồng): \shortans{15\,000}\par\vspace{0.8mm}
\vspace{0.8mm}
\noindent\textbf{Câu 11 (2,0 điểm). Thực hiện phép tính (hợp lý nếu có thể):}
\vspace{0.6mm}

\noindent
\begin{tabularx}{\linewidth}{|X|X|}
\hline
\oBaiTap[9]{a)}{$\dfrac{5}{13} + \left(-\dfrac{8}{3}\right) + \dfrac{8}{13}$} & 
\oBaiTap[9]{b)}{$\dfrac{7}{12} \cdot \left(-\dfrac{11}{19}\right) + \dfrac{7}{12} \cdot \left(-\dfrac{8}{19}\right)$} \\
\hline
\oBaiTap[9]{c)}{$\left(-\dfrac{2}{3}\right)^2 : \dfrac{4}{9} + 4{,}5 \cdot (-2)$} & 
\oBaiTap[9]{d)}{$\sqrt{\dfrac{121}{4}} - |-4{,}5| + \sqrt{9}$} \\
\hline
\end{tabularx}

\newpage
% =============================================================================
% ========================= TOÁN 7 - TUẦN 06 - BUỔI 01: MÃ ĐỀ 104 =============
% ========================= MẶT SAU (TRANG 2/2) ================================
% =============================================================================

\noindent
\begin{tabularx}{\linewidth}{|X|c|}
\hline
\textbf{TOÁN 7 - TUẦN 06 - BUỔI 01} & \textbf{Mã đề: 104 - Trang 2/2} \\
\hline
\end{tabularx}

\vspace{1.2mm}

\noindent\textbf{Câu 12 (2,0 điểm). Tìm số hữu tỉ $x$, biết:}
\vspace{0.6mm}

\noindent
\begin{tabularx}{\linewidth}{|X|X|}
\hline
\oBaiTap[5]{a)}{$x - \dfrac{2}{5} = \dfrac{3}{10}$} & 
\oBaiTap[5]{b)}{$\dfrac{5}{6}x - \dfrac{1}{3} = \dfrac{1}{2}$} \\
\hline
\oBaiTap[5]{c)}{$\left(x - \dfrac{1}{4}\right)^2 = \dfrac{9}{16}$} & 
\oBaiTap[5]{d)}{$|2x - 5| = 3$} \\
\hline
\end{tabularx}

\vspace{1.2mm}

\noindent
\begin{tabularx}{\linewidth}{|X|}
\hline
\vspace{0.5mm}
\textbf{Câu 13 (2,5 điểm).} Cho hai góc kề bù $\widehat{xOy}$ và $\widehat{yOz}$. Biết $\widehat{xOy} = 80^\circ$.\par
\textbf{a) (1,25 điểm)} Tính số đo góc $\widehat{yOz}$.\par
\textbf{b) (1,25 điểm)} Vẽ tia $Om$ là tia phân giác của góc $\widehat{yOz}$. Tính số đo góc $\widehat{xOm}$.\par
\vspace{0.4mm}
\noindent\textit{Bài giải:}\dongke{13}
\\ \hline
\end{tabularx}

\vspace{1.2mm}

\noindent
\begin{tabularx}{\linewidth}{|X|}
\hline
\vspace{0.5mm}
\textbf{Câu 14 (1,0 điểm).} Tìm giá trị nhỏ nhất của biểu thức $A = |x - 7| + 2024$. Chỉ ra giá trị của $x$ để biểu thức đạt giá trị nhỏ nhất đó.\par
\vspace{0.4mm}
\noindent\textit{Bài giải:}\dongke{8}
\\ \hline
\end{tabularx}

\end{document}
